\documentclass[a4paper,oneside,12pt]{extarticle}

\usepackage[utf8]{inputenc}
\usepackage[english]{babel}
\usepackage{amssymb,amsthm,mathtools,enumitem,geometry,hyperref,booktabs}
\usepackage{graphicx}
\usepackage{xcolor}
\usepackage{tikz}
\usetikzlibrary{calc,patterns,decorations.pathreplacing}
\usepackage{pgfplots}
\pgfplotsset{compat=1.18}
\usepackage{float}

\geometry{lmargin=15mm,rmargin=15mm,tmargin=10mm,bmargin=10mm,includeheadfoot}
\pagestyle{plain}
\frenchspacing

\theoremstyle{plain}
\newtheorem{proposition}{Proposition}
\newtheorem{lemma}[proposition]{Lemma}

\theoremstyle{definition}
\newtheorem{definition}[proposition]{Definition}
\newtheorem{example}[proposition]{Example}

\theoremstyle{remark}
\newtheorem{remark}[proposition]{Remark}

\newcommand*{\E}{\mathbb{E}}
\newcommand*{\Var}{\mathrm{Var}}
\newcommand*{\mP}{\mathbb{P}}
\newcommand*{\mR}{\mathbb{R}}
\newcommand*{\ve}{\varepsilon}

\begin{document}
\begin{titlepage}
\begin{center}

{\footnotesize\scshape\color{gray}
Working Notes in Quantitative Risk and Modeling
}

\vfill

{\LARGE\bfseries Homoscedasticity and Its Testing:\\[4pt]
From Conditional Expectation to the Breusch--Pagan Test}

\vspace{36pt}

{\normalsize Boris Garbuzov}

\vfill

{\Large\bfseries Part~4~/~12}

\vspace{6pt}

{\Large\bfseries Fisher Information: Two Routes to the Same Answer}

\vfill

{\small July 2026}

\end{center}
\end{titlepage}

\tableofcontents
\newpage

%====================================================================
\section{Fisher Information for $\mathcal{N}(\mu,1)$}
%====================================================================

\subsection{Definition and operator view}

The \emph{Fisher information} is defined as the variance of the score:
\[
  \mathcal{I}(\theta)
  \;=\; \Var_\theta\bigl[s(X,\theta)\bigr]
  \;=\; \E_\theta\!\left[\left(\frac{\partial}{\partial\theta}
        \log f(X;\theta)\right)^{\!2}\right]
  \;=\; \int_{\mathbb{R}}
        \!\left(\frac{\partial}{\partial\theta}\log f(x;\theta)\right)^{\!2}
        f(x;\theta)\,dx.
\]
The last equality holds in the absolutely continuous case (Lebesgue
density exists).  Since it is an expectation of a square,
$\mathcal{I}(\theta)\ge 0$ always.

Structurally, $\mathcal{I}$ is an operator on one-parametric families
of measures: each family $\{P_\theta\}_{\theta\in\Theta}$ is mapped to
a univariate function $\mathcal{I}(\cdot)\colon\Theta\to[0,\infty)$.
The parameter space $\Theta$ must be a continuum because we need to
differentiate in $\theta$.

The picture below shows the toy case $\Theta=\{\theta_1,\theta_2,\theta_3\}$:
a family of measures on the left maps to three scalar values on the right.

\begin{center}
\includegraphics[width=0.72\textwidth]{nf01_map_discrete.png}
\end{center}

When $\Theta$ is a continuum the output is a curve over $\theta$:

\begin{center}
\includegraphics[width=0.72\textwidth]{nf02_map_cont.png}
\end{center}

%====================================================================
\subsection{Two routes to $\mathcal{I}$}

In the \textbf{generic} (absolutely continuous) case the formula
involves two operations: (1)~square the score, (2)~integrate against
$f_\theta$:
\[
  \mathcal{I}(\theta) \;=\; \int_\mathbb{R} s^2(x,\theta)\,f_\theta(x)\,dx.
\]
Under \textbf{regularity conditions} (differentiation under the
integral sign is valid) a third equivalent formula holds:
\[
  \mathcal{I}(\theta)
  \;=\; -\E_\theta\!\left[\frac{\partial^2}{\partial\theta^2}
        \ln L_x(\theta)\right]
  \;=\; -\E_\theta\!\left[\frac{\partial}{\partial\theta}s_x(\theta)\right].
\]
This uses three operations: (1)~differentiate the score in $\theta$,
(2)~integrate against $f_\theta$, (3)~negate.  We work out both for
$\mathcal{N}(\mu,1)$.

%====================================================================
\subsection{Score surface and the pre-integrand $(x-\mu)^2$}

For the location family $f_\mu(x)=\varphi(x-\mu)$:
\[
  \ell_x(\mu) = -\tfrac{1}{2}\ln(2\pi) - \tfrac{1}{2}(x-\mu)^2,
  \qquad
  s_x(\mu) = \frac{\partial\ell_x}{\partial\mu} = x-\mu.
\]
The pre-integrand for the generic formula is $(x-\mu)^2$ --- a
perfect square.  The bivariate surface $(x-\mu)^2$ over the
$(x,\mu)$-plane:

\begin{center}
\includegraphics[width=0.72\textwidth]{nf03_paraboloid_3d.png}
\end{center}

Fix $\mu=\mu_0$.  The green density $\varphi(x-\mu_0)$ is added
as a section:

\begin{center}
\includegraphics[width=0.72\textwidth]{nf04_paraboloid_sections.png}
\end{center}

%====================================================================
\subsection{Forming the pre-integral curve}

At fixed $\mu_0$, the two factors in the integrand are
$(x-\mu_0)^2$ (green, a parabola) and $\varphi(x-\mu_0)$ (black,
a bell curve):

\begin{center}
\includegraphics[width=0.72\textwidth]{nf05_phi_parabola_2d.png}
\end{center}

Their product $(x-\mu_0)^2\cdot\varphi(x-\mu_0)$ is the
pre-integral curve (red).  It is non-negative and Lebesgue
integrable; its integral is $\mathcal{I}(\mu_0)$:

\begin{center}
\includegraphics[width=0.72\textwidth]{nf06_product_2d.png}
\end{center}

%====================================================================
\subsection{Computing $\mathcal{I}_{\mathcal{N}(\mu,1)}(\mu)$
via the generic route}

Integrating with the substitution $y=x-\mu$:
\[
  \mathcal{I}_{\mathcal{N}(\mu,1)}(\mu)
  = \int_{-\infty}^{\infty}(x-\mu)^2\,\varphi(x-\mu)\,dx
  = \int_{-\infty}^{\infty}y^2\,\varphi(y)\,dy
  = \Var(Z) = 1,\quad Z\sim\mathcal{N}(0,1).
\]
The result does not depend on $\mu$: $\mathcal{I}\equiv 1$ for this family.
In general $\mathcal{I}(\theta)$ depends on $\theta$.

%====================================================================
%====================================================================
\subsection{Regular route: the score surface and its derivative}

The bivariate surface $s(x,\mu)=x-\mu$ over the $(x,\mu)$-plane,
with two sections highlighted:

\begin{center}
\includegraphics[width=0.72\textwidth]{nf07_score_plane.png}
\end{center}

Section at fixed $\mu=x_1$ (score as a function of $x$) ---
a line of slope $+1$ crossing zero at $x=x_1$:

\begin{center}
\includegraphics[width=0.65\textwidth]{nf08_score_line_x.png}
\end{center}

Now differentiate the score in $\mu$:
\[
  \frac{\partial}{\partial\mu}s(x,\mu)
  = \frac{\partial}{\partial\mu}(x-\mu) = -1.
\]
Already at this level the result does not depend on $x$.
The bivariate surface $\frac{\partial}{\partial\mu}s(x,\mu)\equiv -1$
is a flat horizontal plane at height $-1$:

\begin{center}
\includegraphics[width=0.72\textwidth]{nf09_dscore_plane.png}
\end{center}

%====================================================================
\subsection{Sections of the flat plane}

Section at fixed $\mu=\mu_1$ (as a function of $x$) ---
a horizontal line at $-1$, constant in $x$:

\begin{center}
\includegraphics[width=0.65\textwidth]{nf10_dscore_line_x.png}
\end{center}

Section at fixed $x=x_1$ (as a function of $\mu$) ---
also a horizontal line at $-1$, constant in $\mu$:

\begin{center}
\includegraphics[width=0.65\textwidth]{nf11_dscore_line_mu.png}
\end{center}

%====================================================================
\subsection{Product, negation, and the final result}

The product $\frac{\partial}{\partial\mu}s(x,\mu)\cdot f_\mu(x)
= (-1)\cdot\varphi(x-\mu)=-\varphi(x-\mu)$.
At $\mu=\mu_1$ this is the negative density --- an upside-down bell:

\begin{center}
\includegraphics[width=0.65\textwidth]{nf12_neg_phi.png}
\end{center}

Negating gives $\varphi(x-\mu)$, a standard bell curve centred at
$\mu_1$, which integrates to $1$:

\begin{center}
\includegraphics[width=0.65\textwidth]{nf13_pos_phi.png}
\end{center}

\[
  \mathcal{I}(\mu)
  = -\int_{-\infty}^{\infty}
    \frac{\partial}{\partial\mu}s(x,\mu)\cdot\varphi(x-\mu)\,dx
  = \int_{-\infty}^{\infty}\varphi(x-\mu)\,dx = 1.
\]
Both routes confirm $\mathcal{I}_{\mathcal{N}(\mu,1)}\equiv 1$.

\paragraph{Summary.}

\medskip
\begin{center}
\begin{tabular}{lll}
\toprule
Route & Operations & Result \\
\midrule
Generic  & square $s$, integrate            & $\int y^2\varphi(y)\,dy = 1$ \\
Regular  & differentiate $s$, integrate, negate
         & $-\int(-1)\varphi(x-\mu)\,dx = 1$ \\
\bottomrule
\end{tabular}
\end{center}

%====================================================================
\subsection{Existence of Fisher information and the information identity}

\subsubsection*{Existence conditions}

The definition $\mathcal{I}(\theta)=\E_\theta[s(X,\theta)^2]$ presupposes
three things in sequence:
\begin{enumerate}[label=(\arabic*),topsep=2pt,itemsep=1pt]
  \item \textbf{Density exists.}  The family must be absolutely continuous,
        so the likelihood $f(x;\theta)$ (Lebesgue density) is well-defined.
  \item \textbf{Score exists.}  The density must be differentiable in
        $\theta$ (at least a.e.\ in $x$), so that
        $s(x,\theta)=\partial_\theta\ln f(x;\theta)$ makes sense.
  \item \textbf{Score is square-integrable.}  The integral
        $\int s(x,\theta)^2\,f(x;\theta)\,dx$ must be finite.
        In the worst case it diverges: $\mathcal{I}(\theta)=+\infty$.
\end{enumerate}

\subsubsection*{Two functions with the same integral}

Having worked through both routes we can now state the information
identity geometrically.
$\mathcal{I}(\theta)=\E_\theta[s^2]=-\E_\theta[\partial_\theta s]$
says that two \emph{different} functions of $x$ have equal integrals
against $f_\theta$.  For $\mathcal{N}(\mu,1)$ these are
\[
  g_1(x)=(x-\mu)^2\,\varphi(x-\mu),
  \qquad
  g_2(x)=\varphi(x-\mu),
\]
both integrating to $1$.  They look completely different ---
$g_1$ is the two-humped product curve from Section~1.4,
$g_2$ is just the standard bell from Section~1.8 ---
yet their areas coincide.  Left panel: $g_1$ with $\varphi$ and
$(x-\mu)^2$ shown separately; right panel: $g_2=\varphi(x-\mu_1)$
at a fixed $\mu_1$.  Wire-frame and top-down views of the full
bivariate surface $z=(x-\mu)^2\varphi(x-\mu)$ are shown below
the 2D panels:

\begin{center}
\includegraphics[width=0.85\textwidth]{nf14_identity_side_by_side.png}
\end{center}

\subsubsection*{Blue-plane intersection: selecting the 2D curve}

The surface $z=(x-\mu)^2\,\varphi(x-\mu)$ with a vertical blue plane
at $\mu=\mu_0$.  The intersection of the red surface and the blue plane
is exactly $g_1(\,\cdot\,;\mu_0)$; its integral equals $\mathcal{I}=1$:

\begin{center}
\includegraphics[width=0.70\textwidth]{nf16a_z_surface_blue1.png}
\end{center}

A second viewing angle of the same configuration:

\begin{center}
\includegraphics[width=0.70\textwidth]{nf16b_z_surface_blue2.png}
\end{center}

\end{document}
